% ============================================================================
% How to analyse a video and add it to this report; meaning of every export.
% ============================================================================
\section{Procedimiento y exportaciones}\label{sec:procedimiento}

\subsection{Analizar un video}

\begin{enumerate}
  \item \textbf{Abrir} el \texttt{VC\_<código>.MP4} (menú \emph{Archivo}, pestaña \emph{1 · Original}).
  \item \textbf{Escala de tiempo.} En la misma pestaña, dejar marcada \emph{Video acelerado
        (time-lapse)} y fijar $\fr$ (\cref{sec:escala}). Si se conoce la duración real medida,
        usar \emph{Calcular desde la duración real\ldots}.
  \item \textbf{Recorte} (opcional, pestaña \emph{2 · Recortar}): tramo temporal y rectángulo a analizar.
        Los \texttt{VC} ya están recortados a la arena, así que normalmente no hace falta.
  \item \textbf{Seguimiento} (pestaña \emph{4 · Seguimiento}): indicar $N$, verificar el diámetro real
        (\SI{35}{\milli\metre}), indicar si el video está espejado y pulsar \emph{Analizar video}. La calibración del tamaño es automática
        (casilla \emph{Calibrar tamaño automáticamente al analizar}).
  \item \textbf{Revisar} los cuadros de la lista \emph{Cuadros con problemas} (botones \emph{Anterior} y
        \emph{Siguiente}) y, si hace falta, elegir el robot y hacer clic sobre su centro real (ancla). El
        seguimiento se rehace en segundos, sin volver a leer el video.
  \item \textbf{Guardar} el análisis (\emph{Guardar análisis\ldots}, archivo \texttt{.npz}) y exportar el
        CSV con \emph{Exportar CSV (todos los robots)\ldots}, formato \emph{una fila por robot y cuadro}.
\end{enumerate}

\subsection{Generar el informe de ensayos}

Dos scripts de \nolinkurl{software/analisis_video/} arman el informe de resultados
(\nolinkurl{informes/ensayos/informe_ensayos.tex}) a partir de los CSV exportados. Se ejecutan desde
la raíz del repositorio; en Windows la línea se continúa con \verb|^| y en Linux o macOS con
\verb|\|. Por defecto escriben en \nolinkurl{informes/ensayos/generado/} y leen el registro de
presión de \nolinkurl{datos/presion/crudos/}.

\paragraph{Una sección por ensayo.} \texttt{generar\_informe.py} crea
\nolinkurl{generado/<nombre>/} con la tabla de datos, la captura anotada, los gráficos y un
párrafo automático, y actualiza \nolinkurl{generado/lista_ensayos.tex}:
\begin{verbatim}
python software/analisis_video/generar_informe.py ^
       datos/video/trayectorias/VP_20262409_1600_22QR_60_robots.csv ^
       --video datos/video/recortados/VC_20262409_1600_22QR_60.MP4 ^
       --nombre 20262409_1600_22QR_60
\end{verbatim}
Opciones:
\begin{itemize}
  \item \texttt{-{}-cuadros-por-segundo F}: reexpresa el CSV con $\fr = F$. Recalcula $t = n/F$ y
        multiplica velocidades y velocidades angulares por $F/\fr^{\mathrm{CSV}}$, donde
        $\fr^{\mathrm{CSV}}$ se deduce de las columnas \texttt{frame} y \texttt{t\_s}. Corrige un CSV
        exportado con otra escala sin repetir el análisis.
  \item \texttt{-{}-imagen archivo.png -{}-frame n}: usa una imagen ya extraída del cuadro $n$ para la
        captura anotada, en lugar de leer el video (útil si el video no está en la misma máquina).
  \item \texttt{-{}-titulo}, \texttt{-{}-diametro-mm}, \texttt{-{}-diametro-px} (si el CSV no trae escala).
\end{itemize}
Si la duración real supera \SI{10}{\minute}, los ejes de tiempo van en minutos y las trayectorias se
grafican solo para los primeros \SI{5}{\minute}. Las observaciones propias van en
\nolinkurl{generado/<nombre>/notas.tex}, que el script nunca sobrescribe. Si el CSV se exportó con
la corrección de video espejado (columna \texttt{espejo}), la captura se dibuja igual sobre la
imagen original.

\paragraph{Comparación entre ensayos.} \texttt{comparar\_ensayos.py} recibe todos los CSV y genera
\nolinkurl{generado/comparacion/}: tabla comparativa, series temporales con los atascos
marcados, distribuciones, rotación por robot, desplazamiento cuadrático medio, mapas de ocupación y,
si encuentra los registros del sensor, la alineación con el registro del sensor de presión
(\texttt{R\_<código>.csv}) mediante el pulso LED de sincronización y la correlación cruzada:
\begin{verbatim}
python software/analisis_video/comparar_ensayos.py ^
       "datos/video/trayectorias/VP_*_robots.csv"
\end{verbatim}
Después se compila \texttt{informes/ensayos/informe\_ensayos.tex} dos veces (\texttt{pdflatex} o
\texttt{latexmk -pdf}). Este informe (\texttt{informe\_aplicacion.tex}) no depende de los datos.

\subsection{Archivos exportados}

\paragraph{CSV de seguimiento (una fila por robot y cuadro).} Es el archivo de datos del ensayo y la
entrada del informe. Tiempos, velocidades y velocidades angulares están en tiempo real, con la
escala $\fr$ vigente al exportar (\cref{tab:columnas}).

\begin{table}[H]
  \centering
  \caption{Columnas del CSV de seguimiento.}
  \label{tab:columnas}
  \small
  \begin{tabular}{@{}llp{8.6cm}@{}}
    \toprule
    Columna & Unidad & Contenido \\
    \midrule
    \texttt{frame} & --- & número de cuadro $n$ en el archivo de video \\
    \texttt{t\_s} & \si{\second} & tiempo real $n/\fr$ \\
    \texttt{particle} & --- & identificador del robot ($0\ldots N-1$), fijo durante todo el ensayo \\
    \texttt{x\_px}, \texttt{y\_px} & px & centro, relativo al rectángulo analizado ($y$ hacia abajo) \\
    \texttt{status} & --- & estado de la posición (\cref{tab:estados}); \texttt{status\_code}, su código numérico \\
    \texttt{vx\_px\_s}, \texttt{vy\_px\_s} & \si{px\per\second} & componentes de la velocidad \\
    \texttt{speed\_px\_s} & \si{px\per\second} & rapidez $|\vec v|$ \\
    \texttt{vel\_dir\_deg} & \si{\degree} & dirección $\varphi$ de $\vec v$, antihoraria, \SI{0}{\degree} = derecha \\
    \texttt{theta\_deg} & \si{\degree} & orientación acumulada $\theta$, cero en el primer cuadro del tramo \\
    \texttt{omega\_deg\_s}, \texttt{omega\_rad\_s} & \si{\degree\per\second}, \si{\radian\per\second} & velocidad angular $\omega$ \\
    \texttt{rot\_status} & --- & \emph{medida}, \emph{interpolada} o \emph{no disponible} \\
    \texttt{x\_m}, \texttt{y\_m} & \si{\metre} & posición en SI (si se cargó el diámetro real) \\
    \texttt{vx\_m\_s}, \texttt{vy\_m\_s}, \texttt{speed\_m\_s} & \si{\metre\per\second} & velocidad en SI \\
    \texttt{mass}, \texttt{ring\_cov} & --- & intensidad de la detección y cobertura del anillo \\
    \texttt{x\_video\_px}, \texttt{y\_video\_px} & px & centro en coordenadas del video completo (sin corregir espejado) \\
    \texttt{espejo} & --- & solo si se corrigió el espejado: \texttt{horizontal} o \texttt{vertical} \\
    \bottomrule
  \end{tabular}
\end{table}

\paragraph{CSV de seguimiento (una fila por cuadro).} Mismos datos en formato ancho: una fila por
cuadro y, por cada robot $k$, las columnas anteriores con el prefijo \texttt{r}$k$\texttt{\_}
(\texttt{r00\_x\_px}, \texttt{r00\_y\_px}, \ldots). Cómodo para planillas; el informe usa el formato largo.

\paragraph{Resumen por robot.} Una fila por robot: duración (s reales), porcentaje de posiciones
observadas y de rotación medida, rapidez media y máxima, recorrido total, desplazamiento neto, giro
total y $\omega$ media y media en valor absoluto, en píxeles y en SI.

\paragraph{Video con seguimiento.} El tramo analizado con la superposición de la aplicación
(\cref{fig:superposicion}): círculo coloreado por estado, número de robot, flecha de velocidad
(desplazamiento en \num{7.5}~cuadros) y aguja de orientación. Se reproduce a $\fa$, igual que el original.

\paragraph{Sesión de análisis (\texttt{.npz}).} Guarda candidatos, firmas de rotación, anclas
manuales y resultado del seguimiento, para retomar el trabajo sin volver a leer el video. Al
cargarla se aplica la escala de tiempo vigente en la aplicación.
